% NAF 5161, batch 2 — Gallica views 21–40.
% Pass: Opus 5.5 (claude-opus-5-5), 2026-09-22 — first pass, unchecked against the views by a human.
%
% What the views are. Greyscale microfilm scans, portrait, 4239 × about 6120,
% one page per view; leaves of several sizes mounted on guards on a grey card
% ground. Every written view was read from the local mirror as full-width
% bands of the text block at 2400 px, and the fast hand of views 29–30 as a
% grid of three overlapping columns per band at 2000 px (about 2×), with
% tight crops for numbers, figures and doubtful words. \page{N} is always the
% view.
%
% What the twenty views hold. Four pieces, in four hands, none dated.
%
%   v. 21–22  Fol. 10r and its verso, blank. The recto carries only « 10 »
%             (agreeing with the 1888 certificate on v. 2, « Le Feuillet 10
%             est blanc »). Absent from the file.
%   v. 23–26  PIECE 3 (fol. 11r–12v). A short letter-memoir in French, one
%             regular cursive, headed « Qu'il est faux que les Equations qui
%             ne montent que jusques au quarré soient toutes composées de
%             celles dont le Methodiste s'est servy & sa resolution pretendue
%             du Lieu ad quatuor lineas », opening « Monsieur » and addressed to
%             someone who is « ami de M. des C. » and who had asked the writer
%             for a « resolution ». Quotes the Géométrie (« Au reste, dict il, a
%             cause que les equations qui ne montent que jusque au quarré… »,
%             with a marginal « pag. 339 » or so), then constructs a
%             hyperbola as the locus of I with rectangle FB·BI minus rectangle
%             DB·AC equal to a given space X, between asymptotes XN, XS, citing
%             Apollonius (« 4. p. 2 », « 12. prop. 2 ») and naming « le pere
%             Cavalieri » and « Monsieur Mydorge »; then sets up the same locus
%             by algebra: AC = b, FB = c, DB = e, BI = a, whence « ae' + ee
%             sera egal a + ba + x ». One figure (v. 24). It stops at the head
%             of v. 26 after two lines, without closing formula, signature or
%             date; the rest of fol. 12v is blank.
%   v. 27–28  PIECE 4 (fol. 13r–v). A working note on a smaller leaf, in a
%             second, quick, detached hand: to find a primitive right
%             triangle whose area is six times a square and which is not
%             similar to 3, 4, 5 — form a second triangle from the two smaller
%             sides, a third from the two greater, and divide. Worked tables
%             3 4 5 → 7 24 25 → 49 1200 1201, and the quotients 7/10, 120/7,
%             1201/70. The verso (v. 28) has two further tables written
%             sideways (the same chain from 5 12 13: 119 120 169 → 14161
%             40560 42961), a stray « 13 12 5 | 30 », and two pen drawings of
%             tubes or cylinders without lettering. No title, no signature.
%   v. 29–30  PIECE 5 (fol. 14r–v). « Copie de la lettre de Mr Roberval »
%             (title in the hand of the text), in a third, very fast and
%             ligatured cursive, the hardest of the batch: « J'ay 4 choses a
%             repliquer a Mons.r des Cartes ». On the centre of agitation (or
%             of percussion) and the length of the simple pendulum equivalent
%             to a compound one: Descartes (« M. D. C. ») is said to have left
%             out the direction of the motion of the points and to have placed
%             the centre too near the axis; appeals to « l'experience »; ends
%             with two points left to « le R. P. Mersenne ». Signed, large, at
%             the foot of v. 30: « Roberval ». Many words of this hand stay
%             \ill{}; see below. A modern ink note in the left margin of v. 29:
%             « (Du portefeuille 1848 de Libri.) ».
%   v. 31–39  PIECE 6 (fol. 15r–18v). A fair copy in a fourth hand, a round,
%             posed, very regular script with capital letters for the points:
%             an anonymous refutation, addressed to an unnamed « vous », of
%             Descartes's centre of agitation, on the example of a sector of a
%             right cylinder swung about its axis. It opens « Nous conuenions
%             de Definition Monsieur Descartes & moy, touchant le poinct qu'il
%             appelle le Centre d'agitation, lequel nous nommons icy le centre
%             de percussion », states the writer's centres of gravity (O) and
%             of percussion (Q), sets out and criticises the construction of
%             « M. D. C. » (pyramids on cylindrical surfaces, a « triligne aigu
%             parabolique » with centre P at three quarters of the radius),
%             and ends on v. 39 with the experiments that « se sont accordées
%             d'assez pres avec mes conclusions du Centre d'agitation ». The
%             figure is on its own smaller leaf, fol. 17 (v. 36), bound
%             between fol. 16v (v. 35) and fol. 18r (v. 38), which take up the
%             same sentence across it. No signature, no date.
%   v. 34     A second exposure of fol. 16r (the same page as v. 33, line for
%             line). Given as \page{34} with a note, the text not repeated.
%   v. 37     Fol. 17v, blank (the figure shows through). Absent.
%   v. 40     Fol. 19r, blank. Absent. Its verso, v. 41 (batch 3), is an
%             address panel « Pour le Reuerend Pere Mersenne ».
%
% Who writes to whom. Piece 3 names neither writer nor addressee; its
% addressee is a friend of « M. des C. ». Piece 4 is an unsigned note. Piece 5
% is headed as a copy of a letter of Roberval, and signed « Roberval »; it is
% a reply to Descartes and refers the last two points to Mersenne; whether the
% signature is autograph the leaf does not say (its stroke is slower and
% heavier than the text's). Piece 6 names no author; it speaks of Descartes in
% the third person and in the first person of « mes conclusions ». Its last
% paragraph and the fourth point of piece 5 make the same objection in nearly
% the same words (the centre of gravity causes the reciprocation of the swing
% « de droicte à gauche, et de gauche à droicte »); the pass records the
% agreement and attributes nothing.
%
% The hands.
%   Piece 3: a regular, slightly sloping cursive, long filler strokes at line
%   ends (not deletions); final n often a descending loop (« resolutioʒ »),
%   read n; frequent corrections written above a struck word. The letter
%   read D in « D B » (v. 24–26) is a looped capital that could also be B.
%   A small mark stands above « quarré » (title, v. 23) and after « ae » in
%   the equation (v. 24, v. 26); it is kept as an apostrophe in the equation
%   and noted elsewhere.
%   Piece 4: quick, detached letters; Δ for « triangle »; « sex.ple » for
%   sextuple; « R. du q. » for racine du quarré.
%   Piece 5: very fast, ligatured, with dropped endings and long looped
%   descenders; « perpendic. » and a barred p for « per- »; « M. D. C. » for
%   Descartes; « 2e », « 3e », « 4e » with superscript e. Read line by line
%   at 2×; about a third of its words are \ill{} or \uncertain{}.
%   Piece 6: a professional fair copy; u for v inside words and v for u at
%   the head (« conuenions », « vn »), both « et » and « & », « M. D. C. »,
%   « perpend. », « perpendic. », « Exemp. ». Points of the figure are named
%   by letters and by the figures 3 to 9 (« le poinct 3 », « 8 -- 3 », « I 5 »).
% The long s is set s throughout. Spelling as on the leaves (subject,
% connoistre, esgal, aissieu, poinct, droict, touchantes, Appollonius).
%
% Notation. Piece 3 writes its equation with the word « egale » / « sera egal
% a » and lower-case letters a, b, c, e, x; it is set in \[ \] with the words
% in \text{} and the signs as written. Piece 4's tables are set as arrays in
% the page's column order. No other formula in the batch.
%
% Numbers on the leaves. Two series, at the top right of the rectos:
%   (a) « 10 » (v. 21), « 11 » (v. 23), « 12 » (v. 25), « 13 » (v. 27),
%       « 14 » (v. 29), « 15 » (v. 31), « 16 » (v. 33 = v. 34), « 17 »
%       (v. 36), « 18 » (v. 38): one per leaf, the larger figure, the same
%       series as batch 1's 1–9 and in agreement with the 1888 certificate
%       (« Volume de 20 Feuillets / Le Feuillet 10 est blanc »). It behaves as the library's foliation, but it is in
%       ink, like the older series beside it (checked on the images of v. 3 and
%       v. 31), so, as for Latin 17859, NAF 5176 and Français 12357, no
%       \folio{} is written until a human has looked (the coordinator's
%       reconciliation of the batches). On v. 40 (fol. 19r, blank) the corresponding figure
%       is under an ink blot, only a final loop visible; not written.
%   (b) A smaller figure beside or above it, finer pen: « 73 » (v. 25), « 74 »
%       (v. 27), « 48 » struck through with an oblique stroke (v. 29), « 58 »
%       (v. 31), « 59 » (v. 33), « 60 » (v. 36), « 61 » (v. 38), « 62 »
%       (v. 40). On v. 23 a single stroke shaped like a 7 follows « 11 »; it
%       may be the start of a figure of this series (72 would fall between
%       batch 1's 71 and the 73 here), but no second figure shows. The series
%       breaks between fol. 13 (74) and fol. 14 (48), and again between
%       fol. 14 (48) and fol. 15 (58): earlier numberings of leaves that came
%       from different places. Recorded in notes, never as \folio{}.
%   Other numbers: « pag. 339 » (doubtful last figure) in the margin of v. 23,
%   the writer's reference to the page of the work he quotes; the stamp
%   « BIBLIOTHÈQUE NATIONALE MANUSCRITS » (R. F.) on v. 23, 27, 29, 31, 36,
%   39; the numbers inside piece 4's tables and inside the figure of v. 36
%   are content.
%
% Continuity. Piece 3 runs v. 23 → 24 → 25 → 26, each page ending
% mid-sentence. Piece 5 runs v. 29 → 30. Piece 6 runs v. 31 → 32 → 33 → 35
% → (figure leaf v. 36) → 38 → 39. Nothing continues from view 20 into 21
% (fol. 10 is blank) or from view 40 into 41 (fol. 19r is blank; its verso
% is an address panel to Mersenne, belonging to no text in this batch).

\documentclass[11pt,a4paper]{article}
% The preamble shared by every transcription of the mathematicians' manuscripts in Gallica.
%
% It rests on one idea: **uncertainty compounds, it does not resolve.** A
% transcription that smooths over a doubtful word has destroyed the only thing
% it was made for — a reader must be able to tell, at every point, what was on
% the page and what was supplied. The apparatus macros below are therefore the
% critical apparatus, not decoration.
%
% The same commands are recognised by `scripts/render.mjs`, which produces the
% site's reading view, and by `scripts/tei.mjs`. A macro added here must be
% added there too, or rendering fails loudly — which is the wanted behaviour.
%
% Comments here are English, like the rest of the repository. The documents
% this preamble sets are French, and that is deliberate: see the skills.

\ProvidesPackage{germain}[2026/09/15 Transcriptions of mathematicians' manuscripts in Gallica]

\usepackage{amsmath,amssymb,amsthm}
\usepackage{mathrsfs}
% Kept from the parent preamble so the renderer's diagram support stays
% available; these manuscripts draw no commutative diagrams, and nothing here uses it.
\usepackage{tikz-cd}
\usepackage[english,french]{babel}
\usepackage{fontspec}

% --- Greek ------------------------------------------------------------------
% The text face stays fontspec's default, Latin Modern, which has no Greek:
% XeTeX drops every glyph it lacks with a line in the log and nothing on the
% page, and the Greek words the manuscripts quote vanished from the PDFs. So
% the Greek and Greek Extended blocks get a XeTeX character class of their own,
% and entering or leaving a run of it switches the family to CMU Serif — the
% same Computer Modern design, with polytonic Greek — and back. Only the
% family changes: a Greek word in \textit or \textbf keeps its shape and series.
% The transitions to and from every other class carry the tokens that class
% has towards an ordinary letter, so French spacing before « ; » or inside
% « » still holds after a Greek word. No grouping, so no brace or \endgroup
% can come out unbalanced; maths is left alone, since XeTeX inserts nothing
% there. Kept identical in `exercices.sty`.
\makeatletter
\newfontfamily\germainGreekfont{cmun}[
  Extension=.otf, NFSSFamily=germaingreek,
  UprightFont=*rm, ItalicFont=*ti, SlantedFont=*sl,
  BoldFont=*bx, BoldItalicFont=*bi, BoldSlantedFont=*bl]
\newXeTeXintercharclass\germain@greek
\def\germain@greekrange#1#2{%
  \@tempcnta=#1\relax
  \loop
    \XeTeXcharclass\@tempcnta=\germain@greek
    \ifnum\@tempcnta<#2\advance\@tempcnta\@ne\repeat}
\germain@greekrange{"0370}{"03FF}
\germain@greekrange{"1F00}{"1FFF}
\def\germain@greekprev{\rmdefault}
\def\germain@greekin{%
  \edef\germain@greekprev{\f@family}\fontfamily{germaingreek}\selectfont}
\def\germain@greekout{\fontfamily{\germain@greekprev}\selectfont}
\def\germain@greekjoin{%
  \XeTeXinterchartokenstate=\@ne
  \@tempcnta=\z@
  \loop
    \ifnum\@tempcnta=\germain@greek\else
      \XeTeXinterchartoks\germain@greek\@tempcnta=\expandafter{%
        \expandafter\germain@greekout\the\XeTeXinterchartoks\z@\@tempcnta}%
      \XeTeXinterchartoks\@tempcnta\germain@greek=\expandafter{%
        \the\XeTeXinterchartoks\@tempcnta\z@\germain@greekin}%
    \fi
    \ifnum\@tempcnta<4095\advance\@tempcnta\@ne\repeat}
% babel-french sets its own classes, and saves and restores the interchar
% state, each time a language is selected: join again after it does.
\addto\extrasfrench{\germain@greekjoin}
\addto\extrasenglish{\germain@greekjoin}
\AtBeginDocument{\germain@greekjoin}
\makeatother

\usepackage{geometry}
\geometry{a4paper,margin=25mm,marginparwidth=28mm}
\usepackage{marginnote}
\usepackage{xcolor}
\usepackage[normalem]{ulem}
\setlength{\emergencystretch}{3em}
\usepackage{hyperref}
\hypersetup{colorlinks=true,linkcolor=black,urlcolor=black,pdfborder={0 0 0}}

% --- Batch metadata ---------------------------------------------------------
% Taken as they stand by the rendering script, which shows them at the head of
% the reading view. They fix what the file claims to cover. The macro names are
% the parent project's, so that the scripts stay shared; \folder here means the
% volume's slug — `fr-9115` — and \pages counts Gallica views.

\newcommand{\folder}[1]{\gdef\grothFolder{#1}}
\newcommand{\batch}[1]{\gdef\grothBatch{#1}}
\newcommand{\pages}[2]{\gdef\grothFirst{#1}\gdef\grothLast{#2}}
\newcommand{\foldertitle}[1]{\gdef\grothFoldertitle{#1}}

% The holder's own shelfmark, printed as they write it: « Français 9115 ».
\newcommand{\shelfmark}[1]{\gdef\grothShelfmark{#1}}

% The dating, copied verbatim from the holder's catalogue — never inferred
% here. « XIXe siècle » is the BnF's; a pass that dates a leaf from its content
% says so in a \note{}, as its own inference.
\newcommand{\dating}[1]{\gdef\grothDating{#1}}

% The status notice, printed in the title block and repeated as a diagonal
% watermark on every page of the PDF. The manuscripts are in the public
% domain; what this line declares is not a right but a quality — an
% unverified first pass by a machine. The skills write the line; a file
% without it gets no watermark, so leaving it out is visible in the output.
\newcommand{\watermark}[1]{\gdef\grothWatermark{#1}}

\gdef\grothFolder{?}\gdef\grothBatch{}\gdef\grothFirst{?}\gdef\grothLast{?}%
\gdef\grothFoldertitle{}\gdef\grothShelfmark{}\gdef\grothDating{}\gdef\grothWatermark{}

% --- The résumé -------------------------------------------------------------
\newenvironment{resume}
  {\small\setlength{\parskip}{0.45em}}
  {\par\medskip\hrule\medskip}

% --- Keywords ---------------------------------------------------------------
% In English, closing the modernised reading's résumé: the modern vocabulary
% under which to search for what the leaves construct. The manifest extracts
% them to tag the volume — this line is the single source of the tags.
\newcommand{\keywords}[1]{\par\smallskip\noindent
  {\footnotesize\textsc{Keywords} --- \textit{#1}}\par}

% --- The critical apparatus -------------------------------------------------

% The Gallica view number, in the margin. This is the anchor that lets a
% disagreement over a formula be settled by looking at *one* image rather than
% a volume of 750. The site uses it to turn the facsimile as the transcription
% is scrolled. A view is one scanned image: usually one side of a leaf.
\newcommand{\page}[1]{%
  \marginnote{\footnotesize\color{gray!70}\textsf{v.\,#1}}%
  \label{page:#1}\ignorespaces}

% The BnF's pencilled foliation, where the leaf carries it and a pass can read
% it: « 198r ». This is what the literature cites — « ff. 198r–208v » — and
% the one line that turns a citation into a view. Recorded, never inferred:
% a folio a pass cannot read is not written.
\newcommand{\folio}[1]{%
  \marginnote{\footnotesize\color{gray!50}\textsf{f.\,#1}}\ignorespaces}

% The modernised reading's coarser counterpart to \page{}: which views a
% section is drawn from.
\newcommand{\pagerange}[2]{%
  \marginnote{\footnotesize\color{gray!70}\textsf{v.\,#1--#2}}%
  \ignorespaces}

% Illegible: never guessed.
\newcommand{\ill}{\textcolor{red!60!black}{[\,\ldots\,]}}

% A doubtful reading: the word is given, and flagged as uncertain.
\newcommand{\uncertain}[1]{\uline{#1}}

% An editorial addition — a missing letter, an implied word.
\newcommand{\add}[1]{\textcolor{blue!50!black}{[#1]}}

% Struck out by the writer. What was crossed out is part of the document.
\newcommand{\struck}[1]{\sout{#1}}

% The transcriber's note, on the state of the page or the mathematics.
\newcommand{\note}[1]{\par\smallskip{\small\color{gray!60!black}\textit{[#1]}}\par\smallskip}

% A marginal note of hers, distinct from the body of the text.
\newcommand{\marginal}[1]{\par\smallskip\noindent\hspace{1em}%
  {\small\color{gray!60!black}#1}\par\smallskip}

% --- Title ------------------------------------------------------------------
% The title block is French, like the documents themselves.

\AddToHook{shipout/background}{%
  \ifx\grothWatermark\empty\else
    \begin{tikzpicture}[remember picture, overlay]
      \node[rotate=52, scale=3.4, align=center, text=gray!18,
            font=\sffamily\bfseries]
        at (current page.center) {\grothWatermark};
    \end{tikzpicture}%
  \fi}

\AtBeginDocument{%
  \begin{center}
    {\large\bfseries Manuscrits de mathématiciens (Gallica) ---
      \ifx\grothShelfmark\empty\grothFolder\else\grothShelfmark\fi}\\[2pt]
    {\itshape\grothFoldertitle}\\[4pt]
    {\small vues \grothFirst--\grothLast
      \ifx\grothBatch\empty\else\ (lot \grothBatch)\fi}\\[2pt]
    \ifx\grothDating\empty\else
      {\small\color{gray!60!black}Datation du catalogue : \grothDating}\\[2pt]
    \fi
    \ifx\grothWatermark\empty\else
      {\small\bfseries\color{red!45!gray}\grothWatermark}\\[2pt]
    \fi
    {\footnotesize\color{gray!60!black}Transcription établie d'après les images IIIF de Gallica.
     Source gallica.bnf.fr / Bibliothèque nationale de France.\\
     Les lectures incertaines sont soulignées, les illisibles notées [\,\ldots\,],
     les ajouts entre crochets ; \textsf{v.} numérote les vues de Gallica,
     \textsf{f.} la foliotation au crayon de la BnF.}
  \end{center}
  \vspace{1em}}

\endinput


\folder{naf-5161}
\shelfmark{NAF 5161}
\batch{2}
\pages{21}{40}
\dating{1601-1700}
\watermark{Lecture automatique\\non vérifiée}
\foldertitle{Lettres originales de DESCARTES adressées au Père Mersenne (1638-1647), et mémoires divers relatifs aux doctrines de Descartes. II Mémoires divers relatifs aux doctrines de Descartes.}

\begin{document}

\page{23}
\note{En haut à droite, « 11 », de la même main et du même tracé que le « 10 » du feuillet blanc de la vue 21, suivi d'un trait isolé en forme de 7 dont la nature n'est pas claire. Dans la marge de gauche, l'empreinte à l'envers d'un cachet et de quelques lignes : décharge de la page en regard, non transcrite. Le cachet « BIBLIOTHÈQUE NATIONALE MANUSCRITS » (R. F.) est apposé au milieu du texte. Le feuillet est plus court que la page de garde : le bas de la vue est vide.}

Qu'il est faux que les Equations qui ne montent que jusques au \uncertain{quarré} soient toutes composées de celles dont le Methodiste s'est servy \& sa resolution pretendue du Lieu ad quatuor \struck{\ill{}} \uncertain{lineas}
\note{titre en cinq lignes, d'une écriture plus grande et plus appuyée que la suite, chaque ligne prolongée d'un trait de remplissage. « quarré » se termine par une lettre suscrite, peut-être une abréviation (« quarré de quarré » ?) : la lecture n'est pas sûre. Le dernier mot du titre est biffé et récrit au-dessus de la ligne ; la leçon interlinéaire se lit « lineas » ou « lineal ». « le Methodiste » est écrit tel quel.}

Monsieur

Encore que vous soyez ami de M. des C. je croy que vous l'estes assez de la verité pour souffrir que ce qu'il a dit de la composition des lieux solides \struck{\ill{}} imparfaicts, \& defectueux, puis que vous advouez que celuy, dont vous desirez ma resolution ne la peut recevoir par le moyen des equations, dont il faict mention, \& la resolution qu'il pretend donner du lieu des anciens Mathematiciens ad quatuor lineas, sans doute il n'avoit pas \struck{\ill{}} \add{faict} \ill{} \ill{} les denombrements ni les regles generales qui sont la quatriesme partie de son illustre methode, bien que \struck{le} ce grand \& clairvoyant esprit nous proteste d'avoir achevé \& mis la derniere main a tout ce que les autres avoient entrepris \& commencé sur ce subject.
\note{« M. des C. » : abrégé ainsi sur le feuillet. Avant « imparfaicts », un mot bref biffé. « faict » est écrit au-dessus d'un mot biffé de deux lettres. Les deux mots qui suivent « faict » en tête de ligne — le premier commence par la boucle que cette main emploie pour un n final, le second se lit à peu près « indoort » — n'ont pas été lus, même en recadrage serré. « generales » est écrit « g\textsuperscript{e}nales ». « ce » est écrit au-dessus de « le » biffé. La main rend souvent l'n final par une boucle descendante (« resolutio », « moye », « bie », suivis d'une boucle en forme de z) : lu n.}

Au reste, dict il, a cause que les equations qui ne montent que jusque au quarré sont tou\add{tes}
\marginal{pag. \uncertain{339}}
\note{« Au » est écrit au-dessus d'un mot bref biffé. En marge de gauche, en regard de cette ligne, « pag. » suivi de trois chiffres dont le dernier, en forme de q, peut être un 9 : renvoi de l'auteur à la page d'un imprimé, sans doute l'ouvrage qu'il cite ; le chiffre n'est pas sûr. La page s'arrête sur « tou- » ; la suite est à la vue 24. Le « quarré » de cette ligne porte, comme celui du titre, une petite marque suscrite.}

\page{24}
\note{Verso du feuillet de la vue 23 ; la même main poursuit sans interruption. Le feuillet est monté sur un onglet plus grand, dont le bas de la vue montre la partie vide. Aucun numéro. À l'extrême droite de la vue, le bord du feuillet suivant laisse voir quelques lettres qui ne sont pas de cette page.}

toutes composées de celles que je \struck{dont} \add{viens} d'expliquer, non seulement \struck{\ill{}} le probleme des anciens de trois, \& quatre lignes est icy entierement achevé, Mais tout ce qui appartient a ce qu'ils nommoient la composition des lieux solides Et par consequent aussy \&c,
\note{« viens » est écrit au-dessus de « dont » biffé ; en tête de la deuxième ligne, un mot bref biffé. La citation s'arrête sur « \&c ».}

Neantmoins vous souffrirez que si en retranchant la maniere \struck{descrire} \add{de construire} un lieu solide par le moyen de l'algebre on \uncertain{rencontre} \struck{\ill{}} \add{une} equation de mesme forme que celle-cy
\note{« de construire » est écrit au-dessus de « descrire » biffé. « l'algebre » : la fin du mot est surchargée, « re » récrit au-dessus. Le verbe qui suit « on » est lu « rencontre » avec doute ; « une » est écrit au-dessus d'un mot bref biffé. « equation » est abrégé, avec une lettre suscrite.}

\[ ae' + ee \text{ egale } ba + x \]
\note{l'équation est écrite seule au milieu de la ligne. Après « ae », une petite marque suscrite, en forme d'apostrophe, dont la valeur n'est pas claire ; elle est reproduite telle quelle.}

Ceste Equation n'est point composée d'aucunes de celles qu'il s'est donné la peine d'expliquer, Considerez les lignes droites A C, \& B. données de position \& de grandeur, le point B soit a volonté en la ligne droicte \struck{\ill{}} F B prolongée indefiniment \& la ligne B I parallele a C\ill{}. Je dis si le rectangle F B, B I Moins le Rectangle \uncertain{D} B, A C. \struck{est} est egal a l'espace X donné, le point I sera dans une hyperbole, que vous \uncertain{desirez} en ceste sorte. Prenez le point B \struck{\ill{}} \add{ou} vous voudrez a la droicte F B prolongée indefiniment, \& tracez B I parallele a
\note{Entre « Considerez les lignes droites » et le reste du paragraphe, la page porte la figure (voir ci-dessous). « position » est abrégé « po\textsuperscript{on} ». Avant « F B », une lettre surchargée ; après « parallele a C », une lettre noyée sous une tache d'encre. Le premier des deux points du second rectangle se lit mal : D ou B ; la figure porte un point D entre F et B. « desirez » est lu ainsi (d, s long, i) ; le sens appellerait plutôt « descrirez ». « ou » est écrit au-dessus d'une lettre biffée. La page s'arrête sur « parallele a » ; la suite est à la vue 25.}

\note{Figure, à la plume, entre le premier et le second paragraphe. À gauche, un pentagone marqué X (l'espace donné). Au milieu, une droite oblique isolée portant de bas en haut A, C et, au-dessus, E. À droite, une droite horizontale F — D — B, prolongée vers la droite ; de F part une oblique qui passe par x et monte jusqu'à un point marqué n ; par x, une parallèle à F B porte E et S ; de B monte une parallèle à F x qui coupe cette droite en S et s'arrête en I ; une courbe descend de près de n, passe par G (au-dessus de E) et par I, et s'approche de la droite x S vers la droite de la figure.}

\page{25}
\note{Deuxième feuillet de la même pièce, même main. En haut à droite, « 12 », du même tracé que le « 10 » et le « 11 » des vues 21 et 23 ; au-dessus et à gauche, d'une plume plus fine et en plus petit, « 73 », qui paraît appartenir à une numérotation plus ancienne. Dans la marge de gauche, des fragments de lettres et des traits de remplissage : c'est le bord du feuillet précédent, qui dépasse, et non une note de cette page.}

A C, \& de telle grandeur que le Rectangle F B, B I soit egal a l'espace X plus le rectangle A C, \uncertain{D} B, de laquelle B I vous couperez B S egal a A C \& tirez par les points S, F, les deux droictes F X N, S X. celle cy parallele a F B \& celle la parallele a B S \add{\& puis au dedans des asymptotes X N, X S} vous \uncertain{descrirez} une hyperbole par la \uncertain{4}. p. 2 \struck{d'Apollon} d'Appollonius, si vous en desirez connoistre le costé \struck{droit} droict, \& le costé traversant, ou bien en tirant la droicte E G parallele a A C, \& faisant le rectangle X E, E G egal au rectangle X S, \struck{F} S I, car le point G sera un des poincts de l'hyperbole descrite alentour des assymptotes X N, \uncertain{S} X \& qui passe par le point I, \& touts les autres poincts de ceste ligne se trouveront en observant le mesme ordre, ce qui est une maniere aussy courte pour tracer tant des poincts de l'hyperbole qu'on voudra, qu'aucune de celles qui ont esté recueillies par le pere Cavalieri Monsieur Mydorge ou les autres qui ont escrit sur ce subject,
\note{Dans « A C, D B », le premier point du second segment est de nouveau la lettre bouclée qui se lit D ou B (voir la vue 24). « \& puis au dedans des asymptotes X N, X S » est écrit entre les lignes, serré, et rattaché par des renvois (deux petits signes en marge et dans le texte) au début de la ligne « vous descrirez ». Le chiffre qui précède « p. 2 » a la forme d'un 4 ou d'un $\gamma$. Après « d' », « Apollon » est biffé en fin de ligne et le nom récrit en entier à la ligne suivante. « droit » est biffé en fin de ligne et récrit « droict ». Dans « traversant », une lettre est surchargée. Avant « S I », un F biffé. Dans « X N, S X », la première lettre du second nom est la boucle qui, dans cette main, sert aussi bien à S qu'à F ; S est donné d'après « S X » de la ligne 4. « poincts » est écrit avec une surcharge sur la fin du mot.}

au reste je ne m'arreste point a distinguer les differents cas, des diverses solutions ny a la demonstration de ce que j'ay dit cy dessus pource que le tout se deduict facilement de la 12. prop. 2 d'Apollonius. A present il est evident qu'en cherchant ce lieu par l'algebre on trouve une equation de mesme forme que la vostre. car puis que les droites A C, F B sont données nommons l'une b \& l'autre c \& les grandeurs \uncertain{D} B, B I estant vagues \& indeterminées appellons celle cy. e \& celle la a. le rectangle F B, B I
\note{« l'algebre » est de nouveau surchargé en fin de mot. La lettre qui nomme la seconde droite donnée est repassée à l'encre et accompagnée de deux points ; elle se lit c. « indeterminées » : les dernières lettres sont surchargées. La page s'arrête sur « le rectangle F B, B I » ; la suite est à la vue 26. Le bas du feuillet est blanc.}

\page{26}
\note{Verso du feuillet de la vue 25. Deux lignes seulement, en tête, de la même main ; le reste de la page est blanc, l'écriture du recto y transparaissant à l'envers. Aucun numéro.}

sera $ae + ce$. le rectangle \uncertain{D} B, A C $ba$. \& par consequent
\[ ae' + ee \text{ sera egal a } + ba + x \]
\note{« ce » : la première lettre est repassée et plus grasse que les autres ; elle se lit c. Dans la seconde ligne, écrite en retrait comme une formule, « ae » porte la même marque suscrite qu'à la vue 24 ; le signe qui précède « ba » a la forme des autres signes + de la page, avec un crochet à gauche qui le fait ressembler à un 4. La pièce s'arrête ici, sans formule finale ni signature.}

\page{27}
\note{Une autre pièce, sur un feuillet plus petit monté sur onglet, d'une autre main : une écriture rapide, penchée, aux lettres détachées, très différente de celle des vues 23 à 26. En haut à droite, « 13 » (le même tracé que 10, 11 et 12), et à sa droite « 74 », d'une plume plus fine (la série de « 73 » à la vue 25). Le cachet « BIBLIOTHÈQUE NATIONALE MANUSCRITS » (R. F.) est apposé au milieu de la page. Aucun titre, aucune adresse, aucune signature : une note de travail plutôt qu'une lettre.}

Il faut \add{trouver} un $\Delta$ \add{primitif} dont l'aire soit un sex\textsuperscript{ple} quarré et \struck{qui} \add{qui} ne soit \struck{multiple le 3. 4. 5. de 3. 4. 5.} pas semblable au triangle de 3. 4. 5.
\note{« trouver » et « primitif » sont écrits au-dessus de la première ligne, « primitif » rattaché par un renvoi. « sex\textsuperscript{ple} » : l'abréviation de « sextuple », telle qu'elle est écrite. Le triangle est désigné par le signe $\Delta$. Le second « qui » est écrit au-dessus du premier, biffé ; la suite biffée, peu lisible, commence par « multiple » et répète « 3. 4. 5. ».}

car les costez de ce triangle ainsi trouvé estans divisez par la R. du q. dont l'aire du \add{\uncertain{rectangle}} triangle est sextuple l'on aura les \struck{triangle requis} costez du triangle requis.
\note{« R. du q. » : abrégé ainsi (racine du quarré). Le mot écrit entre les lignes et rattaché par un renvoi après « l'aire du » commence par « rect- » et est abrégé ; « rectangle » est une lecture douteuse. « l'on » est écrit en un seul mot et souligné. Un trait horizontal sépare ce paragraphe du suivant.}

pour trouver ce triangle primitif on se servira du triangle donné par exemple 3. 4. 5.

Et des deux moindres costez faire un second triangle

Puis des deux \add{plus} grands costez de ce second triangle \add{faire un 3\textsuperscript{e}} qui sera le requis
\note{« plus » et « faire un 3\textsuperscript{e} » sont écrits au-dessus de la ligne.}

Et \uncertain{dans} ce 3\textsuperscript{e} triangle \add{il \uncertain{convient} \ill{}} que le moindre nombre est impair, et est le quarré du moindre nombre du \struck{pr} second triangle.
\note{Les mots écrits au-dessus de la ligne et rattachés par un renvoi après « triangle » se lisent à peu près « il convient », suivis d'un verbe non lu, de la forme de « trouver » ou « noter ».}

Et le costé pair du mesme 3\textsuperscript{e} triangle est multiple d'un \uncertain{nombre} qu'il y a d'unitez en la moitié de l'aire du 1\textsuperscript{er} triangle.
\marginal{autant}
\note{Dans la marge de gauche, en regard de cette phrase, « autant » précédé d'un signe de renvoi ; un autre renvoi est tracé dans le texte, après « d' » en tête de « d'unitez ». Le mot qui suit « d'un » est en partie sous le cachet.}

\struck{les} costez de ce 3\textsuperscript{e} triangle \struck{estant} divisez par le produit de la racine du costé impair \struck{\ill{}} du mesme multipliée par la \add{racine} \struck{\ill{}} du \struck{\ill{}} du costé pair \struck{du mesme 3\textsuperscript{e} triangle} divisé par le double de l'aire. les trois quotiens seront les costez d'un triangle dont l'aire sera \struck{\ill{}} egale a l'aire du premier triangle.
\marginal{\uncertain{Les}}
\note{Au début de la ligne, « les » est biffé et un mot bref, qui paraît être « Les », est récrit dans la marge. « racine » est écrit au-dessus de la ligne, avec un renvoi. Les mots biffés de ce paragraphe sont brefs et chargés : ils n'ont pas été lus. « 3. 4. » est écrit seul à droite de la dernière ligne.}

\note{Au bas de la page, un tableau de nombres, séparé en colonnes par deux traits verticaux :}
\[
\begin{array}{ccc|c|ccc}
3. & 4 & 5 & & & & \\
7. & 24. & 25 & 3.\ 4 & & & \\
49. & 1200. & 1201. & 24.\ 25. & & & \\
 & & & 49.\ 1200. & \dfrac{7}{10} & \dfrac{120}{7} & \dfrac{1201}{70}.
\end{array}
\]
\note{Au-dessus des trois fractions, une première rangée barrée d'un long trait oblique : le premier terme illisible sur 10, puis $\frac{7}{120}$ et $\frac{70}{1201}$, c'est-à-dire les mêmes fractions renversées. À droite du tableau, isolés et barrés : « 49,12.00 » au-dessus d'un trait, et « 49 00 » au-dessous. À l'extrême gauche du feuillet, au niveau des dernières lignes du texte, trois petits signes superposés (« 1. », « 0 », « 0. ») dont la nature n'est pas claire.}

\page{28}
\note{Verso du feuillet de la vue 27. Pas de texte suivi : deux dessins et trois groupes de calculs, écrits dans des sens différents. Aucun numéro de feuillet. Les chiffres sont de la même plume fine que les calculs de la vue 27.}

\note{Le long du bord droit, écrit perpendiculairement aux lignes (il faut tourner la page d'un quart de tour vers la gauche pour le lire), deux tableaux de nombres côte à côte, chacun fermé à droite par un trait vertical :}
\[
\begin{array}{r|l}
5.\ 12.\ 13 & \\
119..\ 120.\ 169. & 5.\ 12. \\
14161.\ 40560.\ 42961 & 120.\ 169.
\end{array}
\qquad\qquad
\begin{array}{r|l}
3.\ 4.\ 5. & \\
7.\ 24.\ 25. & 3.\ 4. \\
49.\ 1200.\ 1201. & 24.\ 25.
\end{array}
\]
\note{Le premier nombre de la deuxième ligne du premier tableau est mal formé : le chiffre du milieu, courbe, se lit 1 ou 6 ; 119 est donné avec doute. Dans « 42961 », les deux premiers chiffres sont surchargés. Le second tableau reprend celui de la vue 27. Un trait vertical sépare les deux tableaux, et une grosse tache d'encre, étalée, occupe l'espace au-dessous du premier.}

\note{Dans le même sens, plus bas et à gauche de la tache : « 10560 ( 2704. », et au-dessous « 52. » ; à gauche de « 52 », deux ou trois petits chiffres barrés, non lus.}

\note{Au bas de la page, écrit tête-bêche par rapport aux dessins : « 13 12 5 | 30 », les trois côtés du triangle 5, 12, 13 et, après un trait vertical, 30.}

\note{Deux dessins à la plume, sans lettres. À gauche, un tube étroit, ouvert en haut, cerclé d'une bande près de son ouverture, avec une petite pièce en saillie sur le côté. Au milieu, un tube plus large, légèrement évasé vers le bas, dans lequel est logé un cylindre hachuré qui porte à son sommet une petite pièce carrée ; le long du côté droit, une bande étroite en deux segments ; le bas est fermé par une ellipse. Le rapport entre ces dessins et les calculs ne se lit pas sur la page.}

\page{29}
\note{Une nouvelle pièce, d'une quatrième main : une cursive très rapide, serrée, aux mots abrégés et liés, de loin la plus difficile du lot. En haut à droite, « 14 », et à sa droite, barré d'un trait oblique, « 48 », de la série plus ancienne. En marge de gauche, d'une main moderne à l'encre, entre parenthèses : « (Du portefeuille 1848 de Libri.) ». Le cachet « BIBLIOTHÈQUE NATIONALE MANUSCRITS » (R. F.) est apposé dans le bas de la page, sur le texte. Le bas du feuillet est déchiré en festons ; la page est écrite jusqu'au bord. Beaucoup de mots n'ont pu être lus, même à fort grossissement ; seul ce qui a été lu est donné, et chaque lacune est marquée.}

Copie de la lettre de M\textsuperscript{r} \uncertain{Roberval}.
\note{titre, au milieu de la marge supérieure, de la main du texte. Le nom est écrit vite, « Rob » net, la fin en quatre ou cinq jambages : la lecture « Roberval » s'appuie sur la signature de la vue 30.}

J'ay 4 choses a repliquer a Mons\textsuperscript{r} des Cartes. 1\textsuperscript{o}. J'advoue \uncertain{que je me suis} \uncertain{trompé} dans l'escrit que \ill{} \ill{}; mais \uncertain{franchement} \ill{} \ill{} ce que je voiois qu'il fust autheur de la \uncertain{verité} au contraire, je demeure \ill{} \uncertain{qu'aussi tost} qu'il \ill{} \ill{} a \uncertain{ses propositions} il \ill{} \uncertain{necessaire}, et le debat, que s'il estoit capable de la raison et \ill{} \ill{} de Paris \ill{} \ill{}. De quoy je ne puis penser autre chose, sinon que voulant \uncertain{paroistre} \uncertain{infaillible}, il a creu que ce seroit ruiner son dessein, et peut estre aussi qu'il \uncertain{n'auroit} quelque \ill{} \ill{}, s'il se disoit \uncertain{aprentif}. Mais \ill{} \ill{} \ill{}, c'est \ill{} \ill{} l'\uncertain{aiguille}, et \uncertain{de ne vouloir pas} \ill{} \ill{}, \uncertain{s'obstiner} à vouloir maintenir \ill{}, il se \ill{} sans s'en apercevoir, faute que je \uncertain{voy} d'autre \ill{}.
\note{Les six premières lignes. « 1\textsuperscript{o} » est écrit avec un petit o suscrit. La main supprime souvent les finales et lie les mots ; « Mons\textsuperscript{r} » est abrégé ainsi. Entre « le debat » et « De quoy », une demi-ligne n'a été lue que par fragments ; « de Paris » est net.}

\ill{} \ill{} de sa lettre ou il \uncertain{dit} \ill{} \ill{}. Car en sa 2\textsuperscript{e} il n'a que \ill{} \ill{} \ill{} vibration reciproque d'un corps balancé librement à l'entour d'un aissieu, il falloit prendre les distances de \ill{} des parties de ce corps, rapportées à une mesme perpendiculaire \ill{} celle qui \ill{}, \uncertain{dans} cette perpendiculaire le centre d'agitation ou de percussion: et toutesfois dans sa \uncertain{premiere} il avoit asseuré que ce centre estoit dans cette perpendiculaire \ill{}, \uncertain{puisque} selon les regles de la Mechanique, il establissoit \ill{} \ill{} de la force de l'agitation de chaque point du corps balancé, mais aussi de la direction \ill{} de chaque \uncertain{point}, il suppose que la force et la direction \ill{} qui establissent ce centre, doivent estre rapportées à la perpendiculaire dans laquelle il est, et à l'entour de laquelle \ill{} la force de l'agitation \ill{} \ill{}, c'est a dire que les parties distantes ont l'agitation egale a celle des \ill{} \&c. Voila une contradiction manifeste dans sa 2\textsuperscript{e} lettre.
\note{Le paragraphe mêle citation et réfutation ; la ponctuation de la main est rare et celle qui est donnée ici suit les virgules et points visibles. « perpendiculaire » est tantôt écrit en entier, tantôt abrégé par un p barré. Plusieurs mots entre « de sa lettre » et « Car en sa 2\textsuperscript{e} » n'ont pas été lus.}

En 2\textsuperscript{d} lieu \ill{} \ill{} s'il luy plaist, \uncertain{voulant} \ill{} \ill{}, qu'il verra que M. D. C. suppose ce centre dans cette perpendiculaire, et partant estoit obligé de rapporter la force, et la direction de chacun des points agitez (comme \ill{} \ill{}) \ill{} \ill{} \ill{} \uncertain{ce que je luy ay} \uncertain{protesté} 2 fois dans ma \ill{} sur ce point, et (entre autres) quand \ill{} \uncertain{fis voir} que ses raisonnemens estoient \uncertain{defectueux} pour \uncertain{qu'il auroit} oublié de considerer la direction des points mobiles \ill{} qui avoit esté \ill{} \ill{} tous les corps \uncertain{grossiers}, et presque tousjours \ill{}, s'est contenté d'assigner le centre d'agitation ou de percussion plus haut, c'est a dire plus pres de l'aissieu, qu'il n'est en effect. Et quoy que \ill{} \ill{} le lieu ou il est \ill{} cette perpendic. (et non dans les autres) sans \ill{} de \uncertain{m'entendre}, ou de \ill{}; \uncertain{toutesfois} \ill{} \ill{} à \ill{} \ill{}, \ill{} à luy \ill{} \ill{} authorité \ill{} obstiné, ainsi qu'il me reproche sans raison.
\note{Le cachet de la bibliothèque couvre en partie deux lignes de ce paragraphe, vers le milieu de la page.}

La \uncertain{3\textsuperscript{e}} chose que j'en diray avant que \ill{} \ill{} m'a faict \uncertain{connoistre} que ses raisonnemens manquoient en quelque chose, mais je ne suis pas \ill{} \ill{} \ill{} \ill{} au moins, sans \struck{\ill{}} \ill{} prendre \ill{} que \ill{} \ill{}, qui, quoy que j'en die, est trop longue et a laquelle, \ill{} \ill{} authorité je n'ay point pretendu qu'il \uncertain{adjoustast} aucune foy.
\note{Un mot est biffé d'un trait épais au milieu de la deuxième ligne de ce paragraphe.}

En 3\textsuperscript{e} lieu, je n'ay point \ill{} \uncertain{d'indiquer} seulement que le centre de percussion se pourroit assigner \ill{} dans plusieurs autres lignes que dans la perpendic. susdite, dans laquelle \add{seule} il considere \ill{} \ill{} duquel il establit \ill{} regle \ill{}; et que les autres centres estoient dans \ill{} lieu: à quoy M. D. C. \ill{} \ill{} a respondu que ce lieu de \ill{} je \ill{} \ill{} \ill{}.
\note{« seule » est écrit au-dessus de la ligne. Avant « dans plusieurs », une grosse tache d'encre couvre un mot bref. « En 3\textsuperscript{e} lieu » suit, dans la page, un premier « La 3\textsuperscript{e} chose » ; les deux ordinaux sont donnés tels qu'ils se lisent, sans les accorder.}

Pour le 4\textsuperscript{e} je diray plus, car je soustiens que les centres resultans (et non aucun particulier) \ill{} l'effect qu'il attribue en cette occasion au centre de gravité du corps, \ill{} la longueur du pendule qu'il demande. Je soustiens aussi que M. D. C. a manqué de \ill{}, disant que le lieu \uncertain{proposé} ne \ill{} \ill{} en un arc de cercle. Et aussi je soustiens que \ill{} \ill{} ne \uncertain{donne} aucun \ill{} de \ill{} \uncertain{solution}, dont je me suis servi en mes escrits et le tout faute a luy de considerer la direction des points mobiles, laquelle doit necessairement estre considerée en cette occasion.

En 4\textsuperscript{e} et dernier lieu, quand \uncertain{il dit} que quand le centre d'agitation ou de percuss. auroit esté trouvé dans la perpendic., ainsi qu'il pretend, il ne paroissoit pas pourtant qu'il fust la regle ou distance requise pour les vibrations ou balancemens des corps, auquel balancement le centre de gravité contribuë quelque chose aussi bien que le centre d'agitation, veu que ce centre de gravité est la cause de la reciprocation de ce balancement de droite à gauche, et de gauche à droite; et que s'il n'y avoit que l'agitation, le mouvement seroit continuel d'une mesme part à l'entour de l'aissieu; sur quoy M. D. C. dit que c'est la gravité \ill{} \uncertain{supposée} des \ill{} qui est la cause de cette reciprocation de droite à gauche, et non pas le centre de gravité, lequel \ill{} \ill{} \ill{}, \ill{} \ill{} \uncertain{changé} de centre d'agitation.
\note{Ce paragraphe cite, presque mot pour mot, la fin de la copie des vues 38--39 (« il ne paroist pas qu'il fust la regle ou distance requise pour les Vibrations ou balancement des corps… »). La page s'arrête au bas du feuillet, sur un point ; la suite est au verso (vue 30).}

\page{30}
\note{Verso du feuillet 14, même main que la vue 29 ; la lettre s'y poursuit et s'y achève, avec la signature. Aucun numéro. Le haut du feuillet est taché et effrangé. Le texte occupe les trois quarts de la page ; au-dessous, la signature en grand, puis le bas blanc, où transparaît le recto.}

Je \uncertain{ne voy pas} qu'il y ait \uncertain{aucune} \ill{}, et \uncertain{quant} \ill{} il faut la \ill{} et les distinctions qui luy sont attribuées par \ill{} \ill{} dans les \ill{}, et les \ill{}: sans \uncertain{respecter} la Mechanique, dont il \ill{} \ill{} les principes, \ill{} \ill{} \ill{} qu'il \ill{} que quand \ill{} \ill{} \uncertain{pressé} de 2 differentes puissances, chacune a ses \ill{} \uncertain{particuliers}, et alors doit considerer distinctement la force, et la direction de chacune de ces puissances rapportées au centre particulier de chacune; par les moyens de ces differentes forces et directions ainsi rapportées, il establit la force, la direction et le centre du composé de ces puissances; et \ill{} \ill{}, soit que le corps soit libre, ou balancé sur quelque aissieu. Je \uncertain{presume} que les \ill{} considerer distinctement la \uncertain{pesanteur}, et le centre particulier de chacune des parties qui composent un mesme corps, \ill{}, pour les proportions reciproques, establir le centre de gravité du total, lequel centre est le plus souvent tout autre que chacun des particuliers qui pourtant \ill{} \ill{} \ill{}, qui ne se \ill{} pour \ill{} du total.
\note{Les premières lignes, écrites sous la frange tachée du haut du feuillet, sont les plus mal lues de la page.}

\ill{} \ill{}, \ill{} \ill{} l'agitation d'un mesme corps \ill{} d'une autre, quoy qu'elle soit \ill{} de la pesanteur, et chacune de ces puissances a sa force, sa direction, et son centre, \ill{} et particuliers, qui font à examiner le centre du composé de ces differentes puissances, lequel centre change sans doute à chaque differente position du corps balancé à l'entour d'un mesme aissieu; et \uncertain{pour} ce changement, ce centre doit estre \uncertain{cherché} dans le mesme corps. De 2\textsuperscript{e} lieu, les differents poincts estans differemment esloignez de l'aissieu du mouvement, ces differentes distances apportent indubitablement de l'inegalité à la vitesse du mouvement reciproque de ces vibrations. Il faudroit donc \uncertain{examiner}, et en suite \ill{} ces \uncertain{alterations}, pour determiner la longueur du pendule qui feroit les vibrations d'egale durée. C'est de quoy le raisonnement de M. D. C. est par trop esloigné: et partant il ne se faut pas estonner, si l'experience contredit si \uncertain{ouvertement} à sa conclusion, mesme avec corps qui prendent fort peu d'aire, quand \ill{} \ill{} des \uncertain{cercles}, ou plustost des cylindres, balancez à l'entour de leur aissieu, ont \ill{} fort peu d'espaisseur suivant la longueur du mesme aissieu.
\note{« De 2\textsuperscript{e} lieu » : ainsi sur la page. « d'aire » : le mot se lit « daire », d'une seule venue ; la lecture « d'air » serait tout aussi possible.}

Quant au reproche qu'il me fait, que ma \ill{} estoit trop longue: Je n'y respondray autre chose, sinon que sa lettre est aussi trop longue d'autant de mots qu'il y a employé a \ill{} reprocher, qui ne se pratique d'ordinaire, que lors \ill{} \ill{} et sans \ill{} quelque \ill{} de l'amertume, \ill{} manquent d'autres \ill{}, \ill{} à \ill{} de l'amertume. \ill{} \ill{}. Car si M. D. C. \ill{} \ill{} donc à l'\uncertain{une} \add{\ill{} \ill{}} à l'autre, \ill{} aussi, la raison et l'experience sont de \ill{} \ill{}, je laisse à chacun de juger ce que l'on doit \ill{} de luy, qui \uncertain{mesprise} ces 2 instrumens de la connoissance humaine, \ill{} veut suivre que sa pensée, laquelle il \struck{\ill{}} prend \uncertain{luy mesme}, et \ill{} \ill{} passer avec autres, pour autant de \uncertain{verité}.
\note{Au milieu du paragraphe, deux mots sont écrits au-dessus de la ligne, précédés d'un signe de renvoi en forme de double barre, entre « à l'une » et « à l'autre » ; ils n'ont pas été lus. Avant « prend », un mot bref est biffé. Un second signe de renvoi, de même forme, se voit au milieu de la deuxième ligne du paragraphe, avant « quelque ».}

M. D. C. m'a par \ill{} aussi \ill{} \uncertain{Aristarque}. Et dans sa \uncertain{Dioptrique} imprimée, \ill{} \ill{} \ill{} la \ill{} \ill{} \ill{} (ad 3 et 4 lineas) qu'il a faict sur le mesme suiet \uncertain{Apollonius} et Euclide. Mais ces 2 derniers poincts s'adressent au R. P. Mersenne, auquel je \uncertain{satisferay} \ill{}, \ill{} \ill{} chacun \ill{}.
\note{« (ad 3 et 4 lineas) » est écrit entre parenthèses, au milieu de la ligne : c'est le lieu à trois et quatre lignes, qu'examine aussi la pièce des vues 23--26. « Dioptrique » est lu avec doute : le mot commence par une grande boucle et finit par « -ique » ; « imprimée » est net. « R. P. Mersenne » : « R. P. » net, le nom abrégé et lié mais lisible.}

Roberval
\note{signature, au-dessous du texte, à droite, en très grandes lettres à la plume large, R majuscule en boucle, sans paraphe. Le tracé est plus lent et plus appuyé que le texte ; la page ne permet pas de dire si la signature est de la même main que la copie, dont le titre (vue 29) annonce une « Copie de la lettre ». Pas de date, pas d'adresse.}

\page{31}
\note{Une nouvelle pièce, sur un feuillet de grand format, d'une troisième main : une copie au net, en écriture ronde et posée, très régulière, les lettres de la figure en capitales détachées. En haut à droite, « 15 » (le 5 en boucle, du même tracé que 10 à 14), et au-dessus, plus petit, « 58 », de la série plus ancienne. Le cachet « BIBLIOTHÈQUE NATIONALE MANUSCRITS » (R. F.) au milieu de la page. Aucun titre en tête : le texte commence par une grande initiale.}

Nous conuenions de Definition Monsieur Descartes \& moy, touchant le poinct qu'il appelle le Centre d'agitation, lequel nous nommons icy le centre de percussion. Mais sa conclusion est entierement differente de la mienne, de laquelle pourtant, i'ay la demonstration absoluë. Il y a donc quelque deffaut en son raisonnement. C'est ce que ie pretens icy vous faire paroistre. A cét effect, entre plusieurs figures que ie pouuois choisir, je me suis arresté à vn Secteur d'vn Cylindre droict, dans lequel j'espere vous faire voir si clairement ce deffaut qu'il vous sera facile de connoistre qu'il a lieu en toutes les autres figures solides: mesmes en toutes les figures planes desquelles l'aissieu du mouuement n'est pas dans le plan d'icelles, mais perpendiculaire ou oblique à ce plan: et ie croy M. D. C. trop amateur de la verité, pour ne le pas aduoüer, s'il prend la peine de considerer mes raisons.

Soit donc vn Secteur de Cylindre droict A B C D E F G H, duquel l'aissieu tant du Cylindre, que de l'agitation du Secteur, soit la ligne droite A B; ce Secteur estant compris des deux parallelogrammes rectangles A D, A F, qui ont pour costé commun l'aissieu A B; des deux secteurs de Cercles A C G E, B D H F, retranchez des bases du Cylindre; \& de la portion de la Superficie Cylindrique C G E F H D retranchée par ces parallelogrammes et Secteurs de Cercles: et ayant diuisé en deux esgalement l'aissieu A B au poinct I, soit mené par ce poinct vn plan parallele aux bases du Cylindre, lequel plan coupera le Secteur de Cylindre et la Section sera vn Secteur de Cercle, comme I L N M esgal \& parallele aux precedens A C G E, et B D H F: de ce Secteur I L N M, soient les demy-diametres I L, I M; et l'arc L N M, lequel soit coupé en deux esgalement au poinct N auquel soit mené le demy-diametre I N \& prolongé en dehors vers N tant que de besoin. Entendons aussy que cette ligne I N soit perpendiculaire à l'horizon, \& que A B soit de niueau. Dauantage soit I P les trois quarts de I N; et ayant mené L M corde de l'arc L N M, soit entendu que comme l'arc L M est à sa corde L M, ainsy les
\note{La main écrit u pour v à l'intérieur des mots et v pour u à l'initiale (« conuenions », « pouuois », « mouuement », « diuisé », « vn ») ; c'est gardé tel quel. Elle emploie tantôt « et », tantôt la ligature \& : les deux sont gardés. « M. D. C. » : abrégé ainsi sur la page. La phrase se poursuit à la vue 32.}

\page{32}
\note{Verso du feuillet de la vue 31, même main ; la copie se poursuit. Le texte occupe les deux tiers supérieurs de la page ; le bas est blanc. Aucun numéro de feuillet. En haut, au-dessus de la première ligne, de grands traits pâles en forme de lettres : transparence de l'autre face, non transcrite.}

deux tiers du demy-diametre I N, soit à I O portion du mesme demy-diametre: Nous auons demonstré que ce poinct O est le centre de grauité tant du Secteur de Cylindre A H que du secteur de Cercle I L N M. Que si au contraire, on entend que comme la corde L M est a son arc L N M, ainsy soit I P (trois quarts de I N) à I Q portion de I N, nous auons aussy demonstré que le poinct Q sera le centre de percussion ou d'agitation, tant du Secteur de Cylindre A H, que du Secteur de Cercle I L N M.

Toutesfois, suiuant le raisonnement de M. D. C. Il faudroit que ce Centre de percussion ou d'agitation tant du Secteur de Cylindre A H, que du Secteur de Cercle I L N M, fust au poinct P qui est aux trois quarts de la ligne I N; et ce en tout secteur grand ou petit; mesmes au demy-Cylindre, et au demy-Cercle: ce qui est tout contraire à nostre raisonnement, qui fait voir que le veritable centre Q est tousiours plus esloigné de I que P: et ce d'autant plus que le Secteur approchera plus pres d'vn demy-Cylindre, ou d'vn demy-Cercle, n'estant pas toutesfois plus grand: jusques la, que si l'arc estoit d'vn quart plus grand que sa corde, le Centre de percussion du Secteur seroit N: et l'arc estant encore plus grand, ce Centre seroit hors le Secteur, au dela de N.

Mais nostre demonstration est trop longue pour ce lieu; Voyons donc le deffaut de celle de M. D. C. ainsy que nous nous sommes proposez. Et pour ce faire, menons des poincts L, M, les lignes droites L S, M S, qui touchent l'arc L N M, et qui se rencontrent au poinct S dans le demy-diametre I N prolongé. Partant les angles I L S, I M S, seront droicts.

De mesme, ayant pris dans l'arc L N M, deux autres poincts T, V, esgalement esloignez de part et d'autre du poinct N, soient menées les touchantes T R, V R, qui s'entrecoupent au poinct R dans le mesme demy-diametre I N prolongé; et ainsy derechef, ayant mené les demy-diametres I T, I V,
\note{La phrase se poursuit à la vue suivante.}

\page{33}
\note{Recto du feuillet suivant, même main. En haut à droite, « 16 », et au-dessus, plus petit, « 59 ». Dans la marge de gauche, des traces pâles d'écriture : décharge ou transparence, non transcrite. Le bas de la page est blanc.}

les Angles I T R, I V R, seront droicts: Il en sera de mesme de tous les poincts esloignez esgalement de part \& d'autre du poinct N. Enfin, par les lignes A B et I N, soit mené vn plan A B H G, qui coupera le secteur A H en deux autres Secteurs esgaux, et formera le rectangle A B H G, duquel les costez A G et B H couperont aussy en deux esgalement les Secteurs de Cercles A C G E, et B D H F: et par les poincts G, N, H, soient menées des lignes droictes qui touchent les arcs C E, L M, D F, lesquelles touchantes soient Z G 4, X N Y, et 6 H 7, qui seront perpend. aux demy-diametres A G, I N, B H.
\note{Trois des points de la figure sont désignés par des chiffres, 4, 6 et 7, écrits comme les lettres, en capitales détachées ; de même plus bas les points 3 et 8. « perpend. » est abrégé ainsi. La figure elle-même n'est pas sur cette page.}

M. D. Cartes faict donc N X esgale à N Y: puis, dans le demy-diametre ou perpend. I N, ayant pris tel autre poinct qu'on voudra, comme le poinct 3. et par ce poinct entendant vne autre superficie Cylindrique alentour de l'aissieu A B, il veut que comme la Pyramide dont le sommet est I et la base esgale à la Superficie Cylindrique C G H F, est à la Pyramide dont le sommet est I et la base esgale à la Superficie Cylindrique passant par 3. et comprise dans le Secteur A H, ainsy soit l'ordonnée N X a vne autre 3 -- 8 qui luy soit parallele; et ainsy d'vne infinité d'autres poincts que l'on pourra entendre estre trouuez comme ce poinct 8; par tous lesquels poincts vne figure plate estant descrite de part et d'autre de son diametre I N qui la coupe en deux esgalement, Il pretend que le centre de grauité de cette figure platte, sera le centre d'agitation du Secteur A H, ou de tout autre corps pour lequel on aura suiuy les regles de cette construction. Or il est clair que les Pyramides dont il parle, sont icy entre elles comme le quarré de N I au quarré de I 3; \& partant l'ordonnée X N estant à
\note{« 3 -- 8 » : l'ordonnée menée du point 3 au point 8, écrite avec un tiret entre les deux chiffres. « son » dans « de son diametre » est repassé à l'encre. « de I 3 » : le point I suivi du chiffre 3, sans espace. La phrase se poursuit au verso ; la vue 34 reproduit cette même page.}

\page{34}
\note{Seconde prise de vue de la même page que la vue 33 (même feuillet, mêmes numéros « 16 » et « 59 », même texte ligne pour ligne, cadrage et exposition légèrement différents). Le texte n'est pas répété ici ; il est donné à la vue 33.}

\page{35}
\note{Verso du feuillet des vues 33--34, même main. Le texte occupe les deux tiers supérieurs de la page ; le bas est blanc. Aucun numéro de feuillet ; l'écriture du recto transparaît.}

8 -- 3 comme ces Pyramides, c'est à dire, comme le quarré N I au quarré I 3. le centre de grauité de la figure platte (qui est icy vn triligne aigu parabolique) sera au poinct P, qui selon son intention, seroit aussy le centre d'agitation du Secteur A H.
\note{« 8 -- 3 » : l'ordonnée menée du point 8 au point 3, comme à la vue 33.}

Son raisonnement est, que toutes les parties qui sont en la superficie de quelque Cylindre droict duquel A B est l'aissieu, sont esgalement agitées; et que celles qui sont en la Superficie d'vn autre Cylindre plus grand ou plus petit, qui a aussy A B pour aissieu, sont plus ou moins agitées à raison de ce que leur distance de l'aissieu A B est plus ou moins grande. D'ou s'ensuit qu'il y a mesme raison entre la force d'agitation qu'ont ensemble toutes les parties de ce corps, qui sont en la superficie du premier Cylindre, et celle qu'ont toutes les parties du mesme corps, qui sont en la superficie du Second Cylindre qu'il y a entre les Pyramides qui ont leurs bases esgales à ces superficies Cylindriques, et leurs hauteurs esgales aux demy-diametres des mesmes Cylindres. D'ou il suit euidemment (dit il) que le centre de grauité de la figure plate descrite cy dessus, tombe au mesme poinct dans la perpendic. I N, que le centre d'agitation demandé.
\note{« perpendic. » est abrégé ainsi. Dans « qu'ont ensemble », « en » est surchargé.}

Le deffaut de ce raisonnement est, qu'il considere l'agitation seule des parties du corps agité, oubliant la direction de l'agitation de chacune de ces parties; laquelle direction change, et est differente en tous les poincts qui sont inesgalement esloignez du plan vertical A H, quoy que ces poincts soient en vne mesme superficie Cylindrique alentour de l'aissieu A B. Car la direction du poinct L,
\note{La phrase se poursuit à la vue 36 ou au-delà : voir la note de la vue 36.}

\page{36}
\note{Un feuillet plus petit, inséré dans la copie, qui ne porte que la figure. En haut à droite, « 17 », et au-dessus, plus petit, « 60 ». Le cachet « BIBLIOTHÈQUE NATIONALE MANUSCRITS » (R. F.) chevauche le coin supérieur droit de la figure. Le verso (vue 37) est blanc, la figure y transparaissant.}

\note{Figure à la plume, tracée avec soin, en traits pleins et en pointillés, lettres en capitales détachées : le secteur de cylindre droit couché, vu en perspective, son axe A I B horizontal en haut. Aux deux bouts, les secteurs de cercle A C G E et B D H F ; au milieu, la section I L N M, avec l'arc L T N V M renforcé d'un trait gras dans sa partie basse, comme les arcs C G et D H. Sur la verticale I N, prolongée en pointillé jusqu'à S, sont marqués de haut en bas 3, O, P, Q, N, R, 5, puis S. Les touchantes en pointillé : Z G 4 aux arcs de gauche, X N Y au milieu, 6 H 7 à droite ; Des pointillés relient L et M à S, T et V à R, ces points étant sur la verticale. Le point 8 est à gauche de 3, le point 9 à droite, sur la droite pointillée qui va de L vers M. Les chiffres 3 à 9 sont des noms de points, comme dans le texte des vues 33 et 35.}

\page{38}
\note{Recto d'un nouveau feuillet de la copie, même main. En haut à droite, « 18 », et au-dessus, plus petit, « 61 ». Le texte reprend la phrase laissée en suspens à la fin de la vue 35 (« Car la direction du poinct L, ») : le feuillet de la figure (vues 36--37) a été relié entre les deux. Le bas de la page est blanc.}

pour Exemp. est la touchante L S; soit que ce point agité pousse de L vers S, ou qu'au contraire, il tire vers la partie opposée.
\note{« Exemp. » : abrégé ainsi.}

Pareillement, la direction du poinct M, est M S; la direction du poinct T, est T R; la direction du poinct V, est V R; \&c. Tellement que, quoy que l'agitation de tous ces poincts soit esgale, toutesfois la difference de leur direction change l'effect de cette agitation pour deux chefs. Le premier, qu'à l'esgard de la perpendic. I N; ils tirent ou poussent par des poincts differens R, S, \&c. Le Second, que leurs lignes de direction font des angles inesgaux auec cette perpendic. En vn mot, de tous les poincts qui sont en la Superficie Cylindrique C G H F, il n'y a que ceux qui sont dans la ligne G H qui agissent \& fassent leur effort par le poinct N sur la perpendic. I N. tous les autres le faisant en dehors entre N et S: \& partant le centre d'agitation de tous ces poincts, c'est à dire de cette superficie, est aussy entre N et S, et non pas au poinct N, comme il le faudroit pour faire que le raisonnement de M. D. C. fust bon.

De faict, pour auoir ce centre, il faut entendre que comme l'arc L M est à sa corde L M, ainsy le demy-diametre I N soit à I 5, et le poinct 5 sera le centre demandé. Que si on fait le mesme pour toutes les autres Superficies Cylindriques alentour de l'aissieu A B, moindres que C G H F, et comprises dans le secteur A H, on viendra à vne conclusion toute autre que celle de M. D. C.
\note{« I 5 » : le point I suivi du chiffre 5, nom de point de la figure (vue 36), comme 3 et 8 aux vues 33 et 35.}

Je passe soubs silence, que dans toute autre ligne que I N, pourueu qu'elle soit menée du poinct I dans le plan I L N M, on peut assigner vn centre de percussion, et que tous ces centres sont dans vn lieu.

Je passe encor, que quoy que le Centre de percussion ou
\note{La phrase se poursuit au verso (vue 39).}

\page{39}
\note{Verso du feuillet de la vue 38, même main ; la copie s'y achève. Le texte occupe le tiers supérieur de la page ; le reste est blanc, le recto y transparaissant à l'envers (le « 18 » du recto se voit en miroir en haut à gauche). Le cachet « BIBLIOTHÈQUE NATIONALE MANUSCRITS » (R. F.) est apposé sur la dernière ligne. Aucun numéro de feuillet.}

d'agitation fust assigné comme dessus; il ne paroist pas qu'il fust la regle ou distance requise pour les Vibrations ou balancement des corps; auquel balancement le centre de grauité contribuë quelque chose aussy bien que le centre d'agitation. Car ce centre de grauité est la cause de la reciprocation de ce balancement de droicte à gauche, et de gauche à droicte: veu que s'il n'y auoit que l'agitation, le mouuement seroit continuel d'vne mesme part à l'entour de l'aissieu.

Toutesfois, iusques icy, les experiences se sont accordées d'assez pres auec mes conclusions du Centre d'agitation: D'ou i'ay conclu que ce centre d'agitation y contribuë plus que le centre de grauité.
\note{La copie s'arrête ici, au tiers de la page, sans formule finale, sans signature et sans date. Elle ne nomme pas son auteur ; le texte parle de « Monsieur Descartes » et de « M. D. C. » à la troisième personne et s'adresse à un « vous » qu'il ne nomme pas.}

\end{document}
